% 352_Leptonreste_De_ch.tex
% Johann Pascher, 5. September 2026
\providecommand{\BM}{\textbf{[B]}}
\providecommand{\KM}{\textbf{[K]}}
\providecommand{\SM}{\textbf{[S]}}
\providecommand{\QM}{\textbf{[Q]}}

\phantomsection
\addcontentsline{toc}{chapter}{Leptonreste: Leptonleiter, Galois und Koide im Vergleich}

\section*{Statuszeichen}
\begin{center}
\begin{tabular}{@{}lp{10.5cm}@{}}
\textbf{[K]} & \emph{Konfirmiert} --- aus $\xi$ hergeleitet, numerisch geprüft. \\[4pt]
\textbf{[B]} & \emph{Belegt} --- algebraisch bewiesen. \\[4pt]
\textbf{[S]} & \emph{Spekulativ} --- offen. \\[4pt]
\textbf{[Q]} & \emph{Quelle} --- korpusexterne Primärquelle, hier verifiziert. \\
\end{tabular}
\end{center}

\section*{Abstract}

Drei unabhängige FFGFT-Strukturen machen Vorhersagen für Leptonverhältnisse: die
Leptonleiter (Dok.~006/317), die Galois-Identitäten (Dok.~351) und die Koide-Relation.
Dieses Dokument vergleicht alle drei Vorhersagen mit PDG-Werten und drückt die Reste
in Einheiten von $\xi = 4/30000$ aus. Alle Formeln sind exakt rational; numerische
Auswertungen verwenden $m_e = 0{,}51099895$~MeV, $m_\mu = 105{,}6583755$~MeV,
$m_\tau = 1776{,}86$~MeV (PDG) \QM.

\subsection*{Theoretischer Anspruch}

FFGFT beansprucht Vorhersagen für dimensionslose Verhältnisse, nicht für dimensionale
Einzelmassen (Dok.~351, §\,Theoretischer Anspruch). Einzelmassen tragen $K_{\rm frak}$,
SI-Projektion und QED-abhängige Extraktion; diese erzeugen eine unvermeidliche Bandbreite
von Ordnung $\xi$ bis $100\xi$. Die Verhältnisse hingegen sind korrekturfrei ---
\emph{sofern die Korrekturen uniform sind} (vgl.~§\,9).

\section{Massenformel und Quantenzahlen}

Die FFGFT-Massenformel lautet (Dok.~006, 333):
\begin{equation}
  m_i = r_i \cdot \xi^{p_i} \cdot v \cdot K_{\rm frak} \qquad \KM
\end{equation}
mit $\xi = 4/30000$ \BM, $K_{\rm frak} = 74/75$ \KM. Im Verhältnis zweier Leptonen
kürzen sich $v$ und $K_{\rm frak}$ heraus:
\begin{equation}
  \frac{m_j}{m_i} = \frac{r_j}{r_i} \cdot \xi^{p_j - p_i} \qquad \BM
\end{equation}

\begin{center}
\begin{tabular}{@{}lllll@{}}
\toprule
Teilchen & $r_i$ & $p_i$ & $n_\theta,\,n_\phi$ & $r_i = n_\phi^2/n_\theta$ \\
\midrule
Elektron $e$ & $4/3$ & $3/2$ & $3,\,2$ & $4/3$ \\
Myon $\mu$   & $16/5$ & $1$ & $5,\,4$ & $16/5$ \\
Tau $\tau$   & $25/9$ & $2/3$ & $9,\,5$ & $25/9$ \\
\bottomrule
\end{tabular}
\end{center}

\section{Leptonleiter-Vorhersagen}

\begin{equation}
  \frac{m_\mu}{m_e} = \frac{16/5}{4/3} \cdot \xi^{1-3/2}
  = \frac{12}{5} \cdot \xi^{-1/2}
  = \frac{12}{5} \cdot \sqrt{7500}
\end{equation}

\begin{center}
\begin{tabular}{@{}lllll@{}}
\toprule
Verhältnis & Exakte Formel & FFGFT & PDG & Rest \\
\midrule
$m_\mu/m_e$ & $\tfrac{12}{5}\cdot\xi^{-1/2}$ & $207{,}846$ & $206{,}768$ & $+39{,}1\cdot\xi$ \KM \\
$m_\tau/m_\mu$ & $\tfrac{125}{144}\cdot\xi^{-1/3}$ & $16{,}992$ & $16{,}817$ & $+77{,}9\cdot\xi$ \KM \\
$m_\tau/m_e$ & $\tfrac{25}{12}\cdot\xi^{-5/6}$ & $3531{,}6$ & $3477{,}2$ & $+117{,}4\cdot\xi$ \KM \\
\bottomrule
\end{tabular}
\end{center}

Algebraische Abhängigkeit \BM:
\begin{equation}
  \varepsilon(\tau/e) \approx \varepsilon(\tau/\mu) + \varepsilon(\mu/e)
  = 77{,}9 + 39{,}1 = 117{,}0\cdot\xi
\end{equation}
(Der exakte Wert $117{,}4$ enthält einen Kreuzterm der Ordnung $\xi^2$.)
Es gibt nur zwei unabhängige Reste.

\section{Galois-Identitäten}

Aus der Galois-Feldstruktur (Dok.~351) \BM:
\begin{align}
  (m_\mu/m_e)^2 &= 43200 = |GF(9)^*|^2 \cdot 5^2 \cdot |GF(27)| \\
  m_e \cdot m_\mu &= 54~\text{MeV}^2 \quad(\text{Layer-2 mit Anker }m_e)
\end{align}

\begin{center}
\begin{tabular}{@{}lllll@{}}
\toprule
Identität & Galois & PDG & Rest & Status \\
\midrule
$(m_\mu/m_e)^2$ & $43200$ & $42753{,}1$ & $-77{,}6\cdot\xi$ & \KM; Ursache \SM \\
$m_e\cdot m_\mu$~[MeV²] & $54$ & $53{,}991$ & $-1{,}21\cdot\xi$ & emp. Diskrepanz \KM \\
\bottomrule
\end{tabular}
\end{center}

\textbf{Verbindung zur Leptonleiter.} Der Galois-Differenzrest $-77{,}6\cdot\xi$ ist
das Doppelte des linearen Leptonleiter-Rests ($+39{,}1\cdot\xi$) mit umgekehrtem
Vorzeichen --- algebraisch zwingend, da $(m_\mu/m_e)^2_{\rm FFGFT} > (m_\mu/m_e)^2_{\rm PDG}$
genau dann wenn $m_\mu/m_e|_{\rm FFGFT} > m_\mu/m_e|_{\rm PDG}$.

\section{Koide-Relation}

\begin{equation}
  Q \equiv \frac{m_e + m_\mu + m_\tau}{(\sqrt{m_e} + \sqrt{m_\mu} + \sqrt{m_\tau})^2}
  \stackrel{?}{=} \frac{2}{3}
\end{equation}

\begin{center}
\begin{tabular}{@{}llll@{}}
\toprule
Eingabe & $Q$ & Abweichung & Status \\
\midrule
PDG-Massen & $0{,}666\,660\,5$ & $-0{,}046\cdot\xi$ & \KM \\
FFGFT-Vorhersagen (Anker $m_e$) & $0{,}667\,742$ & $+8{,}06\cdot\xi$ & numerisch \\
\bottomrule
\end{tabular}
\end{center}

\textbf{Koide als $m_\tau$-Vorhersage.} Gegeben PDG $m_e$, $m_\mu$ und $Q = 2/3$
folgt algebraisch \KM:
\begin{equation}
  m_{\tau,{\rm Koide}} = 1776{,}969~\text{MeV}
  \quad\text{vs.}\quad
  m_{\tau,{\rm PDG}} = 1776{,}860~\text{MeV}
  \qquad\text{Abweichung: }+0{,}46\cdot\xi
\end{equation}
Das ist die präziseste Vorhersage im gesamten Vergleich.

\section{Gesamtvergleich}

\begin{center}
\begin{tabular}{@{}lllllll@{}}
\toprule
Größe & Methode & Vorhersage & PDG & Rest & Status \\
\midrule
$m_\mu/m_e$ & Leiter & $207{,}846$ & $206{,}768$ & $+39{,}1\cdot\xi$ & \KM \\
$m_\tau/m_\mu$ & Leiter & $16{,}992$ & $16{,}817$ & $+77{,}9\cdot\xi$ & \KM \\
$m_\tau/m_e$ & Leiter & $3531{,}6$ & $3477{,}2$ & $+117{,}4\cdot\xi$ & \KM; Summe \\
$(m_\mu/m_e)^2$ & Galois & $43200$ & $42753$ & $-77{,}6\cdot\xi$ & \KM; Ursache \SM \\
$m_e\cdot m_\mu$~[MeV²] & Galois & $54$ & $53{,}991$ & $-1{,}21\cdot\xi$ & emp. Disk. \KM \\
$Q_{\rm Koide}$ & Koide & $2/3$ & $0{,}66666$ & $-0{,}046\cdot\xi$ & \KM \\
$m_\tau$~[MeV] & Koide & $1776{,}97$ & $1776{,}86$ & $+0{,}46\cdot\xi$ & \KM \\
\bottomrule
\end{tabular}
\end{center}

\section{Befunde}

\textbf{(1) Koide ist am präzisesten \KM.}
PDG-Massen erfüllen Koide auf $0{,}046\cdot\xi$. Koide liefert mit $+0{,}46\cdot\xi$
die genaueste $m_\tau$-Vorhersage im Vergleich. FFGFT-Vorhersagen treffen Koide
dagegen nur auf $+8\cdot\xi$.

\textbf{(2) Leptonleiter und Galois teilen denselben Grundrest \KM.}
Der Galois-Differenzrest ($-77{,}6\cdot\xi$) entspricht exakt $-2\times(+39{,}1\cdot\xi)$
des Leptonleiter-Rests für $m_\mu/m_e$. Dies ist algebraisch zwingend.

\textbf{(3) Alle Reste positiv außer Galois \KM.}
Leptonleiter und Koide-$m_\tau$ liegen knapp über PDG; Galois-$(m_\mu/m_e)^2$ und
$Q_{\rm Koide}$ liegen knapp unter Zielwert. Die Vorzeichen sind konsistent.

\textbf{(4) Bandbreite der Reste \KM.}
\begin{center}
\begin{tabular}{@{}ll@{}}
\toprule
Methode & Rest \\
\midrule
Koide ($Q$) & $\approx 0{,}05\cdot\xi$ \\
Galois (Produkt) & $\approx 1\cdot\xi$ \\
Koide ($m_\tau$) & $\approx 0{,}5\cdot\xi$ \\
Leptonleiter ($m_\mu/m_e$) / Galois (Quadrat) & $\approx 39$--$78\cdot\xi$ \\
\bottomrule
\end{tabular}
\end{center}

\textbf{(5) Koide ist skalenunabhängig \BM.}
Für jeden Faktor $\lambda > 0$ gilt:
\begin{equation}
  Q(\lambda m_1, \lambda m_2, \lambda m_3)
  = \frac{\lambda(m_1+m_2+m_3)}{\lambda(\sqrt{m_1}+\sqrt{m_2}+\sqrt{m_3})^2}
  = Q(m_1, m_2, m_3) \quad\BM
\end{equation}
Konsequenz: Kein uniformer Faktor --- weder $v$, noch $K_{\rm frak}$, noch eine
gemeinsame SI-Projektion --- kann $Q$ verändern. Option~2 (Korrekturen verbessern Koide)
ist damit algebraisch ausgeschlossen.

\section{Skalenunabhängigkeit und ihre Konsequenzen}

\subsection*{Was Koide messen kann}

$Q$ hängt ausschließlich von den Verhältnissen der Massen ab, nicht von deren absolutem
Wert:
\begin{equation}
  Q = Q\!\left(1,\,\frac{m_\mu}{m_e},\,\frac{m_\tau}{m_e}\right)
\end{equation}

Mit den FFGFT-Verhältnissen $(m_\mu/m_e)_{\rm FFGFT} = 207{,}85$ und
$(m_\tau/m_e)_{\rm FFGFT} = 3531{,}6$ ist $Q_{\rm FFGFT} = 0{,}66774$ vollständig
festgelegt --- unabhängig von $v$, $K_{\rm frak}$ oder der SI-Projektion.

\subsection*{Drei getestete Optionen}

\begin{center}
\begin{tabular}{@{}lll@{}}
\toprule
Option & $Q$-Änderung & Status \\
\midrule
Uniformes $K_{\rm frak}$ (alle gleich) & keine & \BM\ ausgeschlossen \\
Einheitlicher $v$-Anker & keine & \BM\ ausgeschlossen \\
Generationsabh.\ $K^{\alpha p_i}_{\rm frak}$ & ja, aber Massen schlechter & Fit ohne Physik \\
\bottomrule
\end{tabular}
\end{center}

\textbf{Generationsabhängiger Test.} Der Exponent $\alpha$ in $m_i\cdot K^{\alpha p_i}_{\rm frak}$,
der $Q = 2/3$ erzwingt, liegt bei $\alpha\approx -36/17$. Die resultierenden Massen
($m_e = 0{,}527$~MeV, $m_\mu = 108{,}0$~MeV, $m_\tau = 1817{,}6$~MeV) liegen weiter
von PDG entfernt als die unkorrigierten Bare-Massen. Keine Physik hinter diesem Wert.

\subsection*{Schlussfolgerung}

Koide ist nicht in der FFGFT-Leptonleiter enthalten \SM. Die Leptonleiter organisiert
die Massen über logarithmische $\xi$-Potenzen: $m_i \propto \xi^{p_i}$. Koide
organisiert die Massen über die Quadratwurzeln $\sqrt{m_i}$. Diese beiden Strukturen
sind orthogonal --- es gibt keine algebraische Beziehung zwischen $\xi^{p_i}$ und
$\sqrt{r_i\,\xi^{p_i}}$ die Koide erzwingt.

PDG-Massen erfüllen Koide auf $0{,}046\cdot\xi$, weil die drei PDG-Massen als
Messsystem zusammenpassen --- sie tragen implizit eine Symmetrie, die in den
FFGFT-Quantenzahlen $(r_i, p_i)$ noch nicht verankert ist.

Der natürliche Kandidat für diese Verankerung ist die $Z_3$-Zirkulant-Struktur des
Leptonmassenoperators (Dok.~A110): dort beschreibt $\theta = 2/9$ die Koide-Phase im
angularen Sektor. Die Frage ist, ob $Q = 2/3$ eine Konsequenz dieser Phase ist.

\section{Koide-Bedingung in den FFGFT-Quantenzahlen}

\subsection*{$Z_3$-Parametrisierung}

Die Eigenwerte eines $Z_3$-Zirkulant-Massenoperators lassen sich schreiben als:
\begin{equation}
  \sqrt{m_k} = c + d\cdot\cos\!\left(\theta + \frac{2\pi k}{3}\right), \quad k = 0,1,2
\end{equation}
Dann gilt $\sum\sqrt{m_k} = 3c$ und $\sum m_k = 3c^2 + \tfrac{3}{2}d^2$, woraus:
\begin{align}
  Q &= \frac{1}{3} + \frac{d^2}{6c^2} \\
  Q = \frac{2}{3} &\iff d = c\sqrt{2} \quad\BM
\end{align}
Das ist die exakte algebraische Bedingung für Koide aus dem $Z_3$-Zirkulant.

\subsection*{$d/c$-Verhältnis in PDG und FFGFT}

\begin{center}
\begin{tabular}{@{}llll@{}}
\toprule
Quelle & $d/c$ & $d/c - \sqrt{2}$ & Status \\
\midrule
PDG & $1{,}41420$ & $-0{,}098\cdot\xi$ & \KM \\
FFGFT-Vorhersagen & $1{,}41649$ & $+17{,}09\cdot\xi$ & numerisch \\
Ziel & $\sqrt{2} = 1{,}41421$ & $0$ & --- \\
\bottomrule
\end{tabular}
\end{center}

PDG erfüllt $d = c\sqrt{2}$ auf $0{,}10\cdot\xi$. FFGFT-Vorhersagen verfehlen es um
$+17\cdot\xi$.

\subsection*{Winkel $\theta$ und Verbindung zu Dok.~A110}

Mit der Zuordnung $k=0\to\tau$, $k=1\to e$, $k=2\to\mu$:

\begin{center}
\begin{tabular}{@{}lll@{}}
\toprule
Quelle & $\theta$~[rad] & $\theta$~[°] \\
\midrule
PDG & $0{,}22223$ & $12{,}733$ \\
FFGFT-Vorhersagen & $0{,}22099$ & $12{,}662$ \\
Dok.~A110 ($\theta = 2/9$) & $0{,}22222$ & $12{,}732$ \\
\bottomrule
\end{tabular}
\end{center}

Alle drei stimmen auf $\lesssim 0{,}5\cdot\xi$ überein \KM. Der Koide-Phasenwinkel
$\theta = 2/9$ aus Dok.~A110 beschreibt die angulare Struktur korrekt. Die Amplitude
$d/c$ hingegen wird nicht durch $\theta$ festgelegt --- sie ist eine unabhängige Bedingung.

\subsection*{Koide-Bedingung in den Quantenzahlen $(r_i, p_i)$}

Da $x_i = \sqrt{r_i}\cdot\xi^{p_i/2}$ und Koide skalenunabhängig ist, hängt $Q$ nur
von den Differenzen der Exponenten ab:
\begin{equation}
  \Delta p_e = p_e - p_\tau = \frac{5}{6},\quad
  \Delta p_\mu = p_\mu - p_\tau = \frac{1}{3},\quad
  \Delta p_\tau = 0
\end{equation}

Die exakte Koide-Bedingung $Q = 2/3$ verlangt:
\begin{equation}
  r_{\mu,{\rm Koide}} = 3{,}2378\ldots \quad\text{statt}\quad r_\mu = \frac{16}{5} = 3{,}2000
\end{equation}
Abweichung: $\Delta r_\mu / r_\mu = +88{,}6\cdot\xi$. Die nächste rationale Näherung
ist $531/164$ --- keine Darstellung als $n_\phi^2/n_\theta$ mit kleinen ganzzahligen
Wicklungszahlen.

\begin{center}
\begin{tabular}{@{}llll@{}}
\toprule
Teilchen & FFGFT $r_i$ & Wicklung $(n_\theta, n_\phi)$ & Koide-$r_i$ \\
\midrule
Elektron & $4/3$ & $(3,\,2)$ & $4/3$ (fest) \\
Myon & $16/5$ & $(5,\,4)$ & $3{,}2378\ldots$ \\
Tau & $25/9$ & $(9,\,5)$ & $25/9$ (fest) \\
\bottomrule
\end{tabular}
\end{center}

Das Koide-$r_\mu$ hat keine saubere Torus-Wicklungsstruktur \SM.

\subsection*{Orthogonalität von $\theta$ und $d/c$ \BM}

Der $Z_3$-Zirkulant hat genau zwei unabhängige reelle Parameter: den Winkel $\theta$
(Rotation im $Z_3$-Raum) und das Verhältnis $d/c$ (Exzentrizität). Ihre Orthogonalität
ist algebraisch bewiesen: $Q = \tfrac{1}{3} + \tfrac{d^2}{6c^2}$ enthält $\theta$
\emph{nicht} --- die Koide-Bedingung ist ausschließlich eine Aussage über $d/c$,
unabhängig von $\theta$ \BM.

Konsequenz für FFGFT: Der Corpus liefert $\theta = 2/9$ aus der Geometrie des
$Z_3$-Zirkulants (Dok.~A110) \KM. Die Bedingung $d/c = \sqrt{2}$ ist eine
\emph{zusätzliche, unabhängige} Forderung --- die R113-Brücke. Beide zusammen
wären nötig, um Koide vollständig aus FFGFT herzuleiten. Aktuell ist nur $\theta$
verankert; $d/c = \sqrt{2}$ bleibt offen \SM.

\subsection*{Schlussfolgerung der Prüfung}

\begin{enumerate}\itemsep4pt
\item \textbf{Winkel $\theta = 2/9$ \KM:} Die angulare Struktur des $Z_3$-Zirkulants
  stimmt zwischen PDG, FFGFT und Dok.~A110 auf $\lesssim 0{,}5\cdot\xi$ überein.
\item \textbf{Amplitude $d/c$ --- zwei verschiedene Objekte:} \KM: Gegeben die
  beobachteten generationsabhängigen $\varepsilon_i$, ist ihre Wirkung auf die
  Koide-Amplitude berechenbar und bringt das bare FFGFT-$d/c$ ($+17\cdot\xi$ neben
  $\sqrt{2}$) quantitativ auf $\sqrt{2}$ --- zu 101\,\% erklärt (s.~§\,9).
  \SM: Die FFGFT-native Herleitung, \emph{warum} die $\varepsilon_i$ diese
  generationsabhängige Struktur besitzen --- und damit warum $d/c = \sqrt{2}$ aus
  der T$^4$/Z$_3$-Geometrie folgen sollte --- bleibt offen (R113-Brücke, §\,10).
\end{enumerate}

\section{Generationsabhängige Korrekturen und Koide}

\subsection*{Die $\varepsilon_i$ sind nicht uniform}

Die bare FFGFT-Massen weichen von PDG-Massen um generationsabhängige Reste $\varepsilon_i$ ab:
\[
  m^{\rm PDG}_i = m^\circ_i \cdot (1 + \varepsilon_i)
\]

\begin{center}
\begin{tabular}{@{}lll@{}}
\toprule
Teilchen & $\varepsilon_i$ & Herkunft \\
\midrule
Elektron & $+192\cdot\xi$ & SI-Projektion, QED \\
Myon & $+152\cdot\xi$ & SI-Projektion, QED-Extraktion (R113) \\
Tau & $+73\cdot\xi$ & SI-Projektion, Messstreuung \\
\bottomrule
\end{tabular}
\end{center}

Die $\varepsilon_i$ fallen mit der Generation --- sie sind nicht uniform.

\subsection*{Warum nicht-uniforme $\varepsilon_i$ Koide verschieben}

Koide ist unter uniformem $\lambda$ skalenunabhängig \BM. Bei verschiedenen
$\varepsilon_i$ gilt in erster Ordnung:
\begin{equation}
  \Delta Q \approx Q^\circ \cdot (\langle\varepsilon\rangle_m - \langle\varepsilon\rangle_{\sqrt{m}})
\end{equation}
Da $\varepsilon_\tau < \varepsilon_\mu < \varepsilon_e$ und $m_\tau$ die Massensumme
dominiert, ist $\langle\varepsilon\rangle_m < \langle\varepsilon\rangle_{\sqrt{m}}$,
also $\Delta Q < 0$.

\subsection*{Numerisches Ergebnis}

\begin{center}
\begin{tabular}{@{}lll@{}}
\toprule
Größe & Wert & Status \\
\midrule
$Q({\rm bare}) - 2/3$ & $+8{,}06\cdot\xi$ & aus Quantenzahlen \\
$\langle\varepsilon\rangle_m$ & $+77{,}7\cdot\xi$ & massegewichtet \\
$\langle\varepsilon\rangle_{\sqrt{m}}$ & $+90{,}0\cdot\xi$ & $\sqrt{m}$-gewichtet \\
$\Delta Q$ & $-8{,}23\cdot\xi$ & Störungsrechnung \\
$Q({\rm bare}) + \Delta Q - 2/3$ & $-0{,}17\cdot\xi$ & Schätzung \\
$Q({\rm PDG}) - 2/3$ & $-0{,}05\cdot\xi$ & exakt gemessen \\
Erklärungsanteil & $101\,\%$ & \KM \\
\bottomrule
\end{tabular}
\end{center}

\subsection*{Revidierte Schlussfolgerung \KM}

Die Formulierung ,,Verhältnisse sind korrekturfrei`` gilt für uniforme Korrekturen \BM.
Für Koide gilt zusätzlich:
\begin{itemize}\itemsep2pt
\item Einfache Verhältnisse $m_j/m_i$ sind unter uniformem $\lambda$ korrekturfrei \BM.
\item Koide $Q$ reagiert auf nicht-uniforme $\varepsilon_i$, weil es über $\sqrt{m_i}$
  definiert ist.
\item Die generationsabhängigen $\varepsilon_i$ (Ordnung $73$--$192\cdot\xi$, fallend
  mit der Generation) erklären den Übergang $Q_{\rm bare}(+8\cdot\xi)\to Q_{\rm PDG}(-0{,}05\cdot\xi)$
  zu $101\,\%$ \KM.
\item Der verbleibende offene Punkt ist: \emph{warum fallen $\varepsilon_i$ mit der
  Generation?} Die Antwort liegt in der QED-Extraktion (R113) und der SI-Projektion
  (Dok.~085, 333) --- deren generationsabhängige Struktur ist in §\,10 analysiert.
\end{itemize}

\section{Warum fallen $\varepsilon_i$ mit der Generation?}
\label{sec:eps-herkunft}

\noindent\emph{Die Statustrennung in diesem Abschnitt wurde in Korrespondenz mit
Stefaan Vossen unter Anwendung des Constitutional Physics / Canonical Constitutional
Record-Rahmens entwickelt (Dok.~351, 5.--6.~September~2026): die berechenbare Wirkung
der beobachteten $\varepsilon_i$ \KM\ wird von ihrer quellnativen Herleitung \SM\
unterschieden.}

\subsection*{Zwei Quellen der nicht-uniformen Korrekturen}

\paragraph{1.\ QED-Selbstenergie (R113) \QM.}
Die gemessene Pol-Masse eines geladenen Leptons unterscheidet sich von der nackten
Masse durch die QED-Selbstenergie. In führender Ordnung gilt:
\begin{equation}
  \frac{\delta m_i}{m_i} \approx \frac{3\alpha}{4\pi}
  \ln\!\left(\frac{m_i}{\mu}\right) + \text{const},
  \label{eq:qed-self}
\end{equation}
wobei $\mu$ die Renormierungsskala ist. Da $m_\tau \gg m_\mu \gg m_e$,
wächst $\ln(m_i/\mu)$ mit der Generation.

Der beobachtete Trend $\varepsilon_e > \varepsilon_\mu > \varepsilon_\tau$ ist jedoch
\emph{entgegengesetzt} zur naiven QED-Selbstenergie, die mit der Masse zunimmt.
Der Grund: $\varepsilon_i$ misst den Abstand bare FFGFT-Masse $\to$ PDG-Masse.
Dieser Abstand ist eine \emph{Differenz} zweier Größen, von denen jede
massenabhängig skaliert. Die SI-Projektion dominiert bei kleinen Massen stärker
(relativer Korrekturbetrag $\propto 1/m_i$ für feste absolute Projektionsunschärfe).

\paragraph{2.\ SI-Projektion (Dok.~085, 333) \SM.}
Die SI-Projektion überführt die bare FFGFT-Masse $m^\circ_i = r_i\cdot\xi^{p_i}\cdot v\cdot K_{\rm frak}$
in den gemessenen MeV-Wert. Der Faktor $K_{\rm frak} = 74/75$ ist uniform \BM\ und
verschiebt $Q$ nicht (§\,5). Die \emph{nicht-uniforme} Komponente entsteht aus der
massenabhängigen QED-Extraktion: bei der experimentellen Bestimmung von $m_\mu$ aus
dem Myonen-Zerfall und von $m_\tau$ aus $\tau$-Zerfallsbreiten gehen jeweils
verschiedene Ordnungen der Störungstheorie ein (R113). Diese sind von Ordnung
$\alpha\approx 1/137$ und damit von Ordnung $\xi$ bis $100\xi$ --- konsistent mit
den beobachteten $\varepsilon_i$.

\subsection*{Was FFGFT herleiten kann und was nicht}

\begin{center}
\begin{tabular}{@{}lll@{}}
\toprule
Größe & Status & Herkunft \\
\midrule
Bare-Massen $m^\circ_i = r_i\cdot\xi^{p_i}\cdot v\cdot K_{\rm frak}$ & \KM & FFGFT intern \\
Uniformer Faktor $K_{\rm frak}$ & \BM & FFGFT intern \\
Größenordnung der $\varepsilon_i$ ($\sim\alpha/\pi$, Ordnung $\xi$--$100\xi$) & \KM & QED extern \QM \\
Fallende Struktur $\varepsilon_e > \varepsilon_\mu > \varepsilon_\tau$ & \SM & QED + SI-Projektion,
  nicht aus FFGFT-Quantenzahlen hergeleitet \\
QED-Selbstenergie aus T$^4$/Z$_3$-Geometrie & \SM & offene Brücke R113 \\
\bottomrule
\end{tabular}
\end{center}

\subsection*{Verbindung zu Dok.~351}

Dieselbe nicht-uniforme $\varepsilon_i$-Struktur erklärt die Verhältnisabweichungen
in Dok.~351: $m_\mu/m_e$ weicht um $\sim 39\cdot\xi$ und $m_\tau/m_\mu$ um $\sim 78\cdot\xi$
von den FFGFT-Vorhersagen ab. Einfache Verhältnisse sind unter \emph{uniformem}
$\lambda$ korrekturfrei \BM; unter nicht-uniformem $\varepsilon_i$ tragen sie
den Differenzrest $\varepsilon_j - \varepsilon_i$. Die Brücke R113 (QED-Extraktion,
massenabhängig) ist damit der gemeinsame offene Punkt von Dok.~351 und Dok.~352
(Dok.~351, §\,5: ,,R113: $m_\mu/m_e$ QED-abhängig``, Textkorrektur ,,modellunabhängig''
$\to$ ,,modellärmst'').

\subsection*{Schlussfolgerung}

Die beobachtete fallende $\varepsilon_i$-Struktur ist mit einer kombinierten
QED/SI-Projektions-Beschreibung verträglich, ihre Herleitung aus den
FFGFT-Quantenzahlen bleibt jedoch offen \SM. FFGFT liefert die Bare-Massen; die
Spreizung der $\varepsilon_i$ ist eine externe \QM-Tatsache in ihrer Größenordnung.
Ob die QED-Selbstenergie aus der T$^4$/Z$_3$-Geometrie herleitbar ist, bleibt offen
\SM\ --- das ist die R113-Brücke.

\bigskip\noindent
\textbf{Prüfskript:}
\texttt{pruef\_352\_leptonreste.py} --- alle Zahlen reproduzierbar, reine Python-Standardbibliothek.

\section*{Quellen}

PDG 2024 \QM. Dok.~006 (Massenformel), Dok.~317 (Leptonleiter-Zahlenwerte),
Dok.~351 (Galois-Identitäten, $\varepsilon_i$-Werte, R113), Dok.~A110
($Z_3$-Zirkulant), Dok.~085, 333 (SI-Projektion), R113 (QED-Abhängigkeit
$m_\mu/m_e$), Dok.~190/R114 (Registereintrag dieses Dokuments).
